\documentclass[10pt]{article}
\usepackage[margin=20mm]{geometry}
\usepackage{amsmath,amssymb,mathtools,fontspec,unicode-math,xurl,hyperref}
\IfFontExistsTF{DejaVu Serif}{\setmainfont{DejaVu Serif}}{\setmainfont{Latin Modern Roman}}
\setmathfont{latinmodern-math.otf}
\setlength{\parindent}{0pt}\setlength{\parskip}{6pt}\sloppy
\title{Technical Verification of
\(C^*\)-Rigidity of Bounded Geometry Metric Spaces}
\author{AI-assisted technical review}
\date{}
\begin{document}
\maketitle
Abstract: The exact publisher version is reviewed through the concentration estimate, the construction of coarse inverse maps, and the refinement giving the outer automorphism group. No substantive mathematical defect is established in the declared inspected scope. The signed-series argument is justified for its actual application, the approximation constants can be chosen strictly, and the strong-limit step follows from weak compactness of the error ball. Local notation and coefficient-space conventions are clarified without converting them into unsupported findings. The identified external uniform-approximation and multiplier inputs are not all independently reproved. Scientific adjudication remains pending.\par
\section*{Paper identity, version and scope}
The paper is Diego Martínez and Federico Vigolo, \(C^*\)-rigidity of bounded geometry metric spaces, Publications Mathématiques de l’IHÉS 141 (2025), 333–348, DOI 10.1007/s10240-025-00155-3 (https://doi.org/10.1007/s10240-025-00155-3) [paper]. This report concerns the acquired 16-page publisher version of record, published online on 4 March 2025. The separately acquired arXiv:2501.03128v1 author source is not asserted equivalent to that version. A physical PDF page is one-based; its printed number is larger by 332.\par
The root reviewer read the complete extracted canonical text. The principal checks concern definitions and quantifiers, Proposition 3.2, Theorems A and B, Proposition 4.3 and Theorem 4.5. Full text reading is distinguished from independent reproof of every external dependency. No separate supplement or original computational repository was identified in the bounded source checks; no defect is inferred from their absence. There are no manufactured numbered Problems.\par
\section*{Operative definitions and external interfaces}
Throughout the operative arguments, \(X,Y\) are uniformly locally finite ordinary metric spaces. Thus nonempty bounded balls are finite, and each nonempty space is countable by exhaustion about a fixed point. The main Roe and controlled-propagation algebras use \(H=\ell^2(\mathbb N)\); the uniform Roe algebra uses \(H=\mathbb C\). The zero Hilbert space is not an alternative coefficient choice for the main rigidity assertion.\par
Write \(P_A\) for the projection onto \(\ell^2(A;H)\). Weak approximate control means that for every propagation bound \(r\) and error \(\varepsilon>0\), one propagation bound \(R\) works for all contractions of propagation at most \(r\). In particular it controls the projections \(P_A\), even when these are not locally compact for infinite-dimensional \(H\).\par
The spatiality interface in Theorem 2.14 can be recovered from the compact operators: they form the unique minimal nonzero ideal in either algebra, and an isomorphism on that ideal is spatial. Multiplying by rank-one operators then shows that the same unitary implements the entire isomorphism. This argument does not require separability.\par
The exact author version of Braga–Vignati [bv] was consulted for the uniform-approximation and multiplier interfaces. Its Theorems 3.4–3.5 do not require property A, despite the title of that paper. For the Roe interface, control initially stated on the Roe algebra extends to every bounded finite-propagation contraction: compress by diagonal finite-rank finite-support projections, use the uniform bound, and take a weak-operator limit of bounded approximants. Propagation constraints and the closed error ball persist under that limit. The bounded auxiliary notes also supply a separable reducing-space argument for arbitrary \(H\); the main coefficient choices already lie within the directly consulted scope.\par
\section*{The concentration estimate and coarse inverse maps}
Equation reference: eq:energy, eq:corner. Consider Proposition 3.2, PDF pp. 7–8 (printed pp. 339–340). Let \(U:\ell^2(X;H)\to\ell^2(Y;H)\) be unitary, choose a unit vector \(h\in H\), and set \[v=U^*(\delta_y\otimes h),\qquad B=B_Y(y;R),\qquad
 w_x=(I-P_B)UP_xv .\] Under its hypothesis \(\|P_BU P_x\|\le\delta\), \[
 \sum_x\|w_x\|^2
 \ge (1-\delta^2)\sum_x\|P_xv\|^2
 =1-\delta^2.\] The original family \(P_xv\) is orthogonal. Its signed sums converge unconditionally, uniformly in the signs on tails. Applying the bounded operator \((I-P_B)U\) proves the same convergence needed for \(w_x\). Finite Rademacher averaging followed by this limit therefore gives \[\mathbb E\Bigl\|\sum_x\epsilon_xw_x\Bigr\|^2
       =\sum_x\|w_x\|^2 .\] Moreover \(\sum_xw_x=(I-P_B)Uv=0\). With \(A=\{x:\epsilon_x=1\}\), the signed sum equals \(2\sum_{x\in A}w_x\). For any \(0<\varepsilon<\frac12\sqrt{1-\delta^2}\), choose \(\eta<1-\delta^2-4\varepsilon^2\). The averaging argument then supplies signs for which \[
 \|P_{Y\setminus B}UP_AU^*P_y\|
 \ge \Bigl\|\sum_{x\in A}w_x\Bigr\|
 >\varepsilon.\] Local finiteness of \(Y\) ensures \(d_Y(y,Y\setminus B)>R\), so [eq:corner] (\#eq:corner) contradicts \(\varepsilon\)-\(R\)-quasi-locality. This verifies the constant \(1/2\), the quantifiers and the strict error margin.\par
The general wording of Lemma 3.1 needs the usual almost-everywhere or finite-partial-sum interpretation: square-summable norms alone do not make every signed series converge, as \(v_n=h/n\) illustrates. Its actual application above has the stronger unconditional convergence, so that wording does not leave a gap in Proposition 3.2. Likewise, in Lemma 2.13 one can choose witness vectors with norms strictly greater than half the asserted lower bound, or use error \(\delta^2/8\). The printed weak inequalities do not obstruct the required strict contradiction.\par
For Theorem A, weak approximate control gives a single \(R\) for all \(UP_AU^*\). The contrapositive of Proposition 3.2 then chooses \(g(y)\) with \(\|P_{B(y;R)}UP_{g(y)}\|>\delta\). Apply Lemma 2.13 to \(U^*\), which implements the inverse isomorphism, and to \(B(y;R)\cup B(y';R)\). Its diameter is at most \(2R+d_Y(y,y')\), giving uniform control of \(g\). The symmetric construction gives \(f\).\par
To check \(g\circ f\), use \(U^*\), the domain ball \(B(f(x);R)\), and output sets \(\{g(f(x))\}\) and \(B(x;R)\). The two norm bounds in the proof give a uniform bound on the distance between these sets; adding \(R\) bounds \(d_X(g(f(x)),x)\). The other composition is symmetric. This reconstruction retains the uniform constants and does not interchange pointwise and uniform choices.\par
\section*{The refinement and the outer automorphism group}
Theorem 4.5, PDF pp. 10–12, assumes infinite-dimensional \(H\). For a finite-rank diagonal projection \(p=\sum_i p_{x_i}\), Claim 4.6 chooses fiber unitaries so that the unwanted output tails have orthogonal ranges as well as orthogonal initial supports. Only finitely many finite-dimensional constraints occur at each step; infinite-dimensionality supplies the needed room. The norm of the sum of those tails is therefore their maximum norm, rather than a bound depending on the number of fibers.\par
The product-approximation step can require contractive approximants without changing propagation: first approximate a contraction with error \(\varepsilon/4\), then rescale by \(\max(1,\|s\|)\). The resulting error is at most \(\varepsilon/2\). This justifies the stated product error with a uniform support radius.\par
Here is a direct justification of the final closure step, also valid for nets. Let \(S_R\) be the linear subspace of operators supported within radius \(R\) of \(f\). Its defining zero matrix blocks make it weak-operator closed. If \[T_\lambda=S_\lambda+E_\lambda,\qquad
 S_\lambda\in S_R,\quad \|E_\lambda\|\le\varepsilon,\quad
 T_\lambda\longrightarrow T\ \hbox{strongly},\] take a weak-operator convergent subnet of the errors in their compact ball, with limit \(E\). Then \(S_\lambda\to T-E\) weakly, so \(T-E\in S_R\) and \(\|E\|\le\varepsilon\). Thus the fixed-radius error neighborhood is strongly closed. No assertion that every strongly convergent net is uniformly norm bounded is needed.\par
Proposition 4.3 preserves local compactness using finite-set tails. For finite \(A\), Lemma 2.13 gives finite neighborhoods \(N_R(B)\) such that \(\|(I-P_{N_R(B)})U^*P_A\|\to0\). For every locally compact \(t\) on the whole space, \(tP_{N_R(B)}\) is compact because that neighborhood is finite. Hence \(UtU^*P_A\) is a norm limit of compact operators; the left condition follows similarly. This supplies the argument with the correct whole-space operator domain.\par
For Theorem B, let \(T_n\to U\) in norm with every \(T_n\) coarsely supported on \(f\), and choose a unitary cover \(W\) of \(f\). Each \(T_nW^*\) has finite propagation: any nonzero block factors through an input point whose two output points are within the two support radii of its image under \(f\). Consequently \[UW^*=\lim_nT_nW^*\in C^*_{\mathrm{cp}}(X)\] is unitary. This proves surjectivity modulo multiplier innerness.\par
Injectivity also has a short direct check. If the cover \(U_f\) induces an inner automorphism, it differs from a unitary in \(C^*_{\mathrm{cp}}(X)\) by a scalar, since the commutant of all compact operators is scalar. Choose finite-propagation \(S\) with \(\|U_f-S\|<1\). If \(d(f(x),x)\) exceeded the sum of the support radius of \(U_f\) and \(\operatorname{Prop}(S)\), the images of the \(x\)-fiber under \(U_f\) and \(S\) would be orthogonal. This would force \(\|(U_f-S)P_x\|\ge1\), a contradiction. Thus \(f\) is close to the identity. Covers of close maps and compositions obey the same support estimates, giving the stated canonical homomorphism.\par
\section*{Clarifications, limitations and conclusion}
The proof uses \(Y\) for output complements and \(B(\ell^2(X;H),\ell^2(Y;H))\) for operators between the two spaces. The local \(X/Y\), tensor-order and whole-space/domain slips in the printed formulas are uniquely resolved by those types and do not invalidate the reconstructed estimates. Implementing unitaries are understood modulo their scalar ambiguity, or as normalizers before taking that quotient.\par
The countable basis in the displayed covering-unitary construction is appropriate to the main coefficient \(H=\ell^2(\mathbb N)\). For other infinite Hilbert dimensions, use the actual basis index set; multiplication by a finite nonzero cell cardinality leaves that infinite cardinality unchanged. Empty-space wording in the elementary identity and finite-dimension remarks is a convention boundary. For the main nonzero coefficients, an isomorphism with one empty space forces both spaces to be empty, where the conclusion is immediate.\par
The review retains the externally proved uniform-approximation and multiplier inputs at their identified interfaces. It does not claim to reprove all of their foundations or the broader arbitrary-proper-space results reserved for another paper. No empirical or author-code verification is relevant to these exact arguments.\par
Dimension | Recorded state\par
Issue assessment | No Issue Found within the declared inspected scope.\par
Numbered substantive problems | None established.\par
Scope clarifications | Convergence interpretation, strict tolerances, operator domains and coefficient conventions checked as above.\par
External proof coverage | Specific interfaces consulted; full foundational reproof not claimed.\par
Scientific adjudication | Pending; no human review supplied.\par
Author adoption / acceptance | No correction proposed for adoption; independent acceptance not performed.\par
9 D. Martínez and F. Vigolo, \(C^*\)-rigidity of bounded geometry metric spaces, Publ. Math. IHÉS 141 (2025), 333–348. doi:10.1007/s10240-025-00155-3 (https://doi.org/10.1007/s10240-025-00155-3). Reviewed publisher version. B. M. Braga and A. Vignati, A Gelfand-type duality for coarse metric spaces with property A, arXiv:2007.10499v2 (https://arxiv.org/abs/2007.10499v2), 28 July 2020. Exact 34-page author version consulted at the stated interfaces; equivalence to the later IMRN version is not asserted.\par
\section*{Supporting checks}
The canonical source PDF SHA256 is \[\texttt{072d5ed20d6afdc9a00fb8ca6befdc102}\] \[\texttt{461adbfef4f4c6f9d42ab81dfed6d19}.\] Source/manifest hashes, the distinct author-source version, proof checkpoints, external interface files, actual source-image observations, compiler/support dependencies and actual report-page checks are retained in the matching verification directory. The arguments are analytical; no author program or author-source compilation was used.\par
The root role read and integrated the complete paper and wrote this report. They are neither human reviewers nor independent acceptance. Scientific decision remains pending and the human-review list is empty.\par
\end{document}
